% TOP_PROBLEM: {"id": 2532, "title": "Kourovka Notebook Problem 21.23", "category_id": 17, "category": {"id": 17, "name": "group_theory", "display_name": "Group Theory", "description": "Problems about groups, group actions, representations, and related algebraic structures.", "slug": "group-theory", "order_index": 17, "created_at": "2026-07-31T15:26:25.671Z"}, "status": "open", "problem_number": "KOU-21.23", "set_id": 10}
% FIRST_POSTED: 2026-10-01
% ENTRY_KIND: statement_only
% UnsolvedMath ID 2532; source catalogue code KOU-21.23
% !TeX program = lualatex
% Definition and sources only; this file is not an LLM solution attempt.
\documentclass[11pt,a4paper]{article}
\usepackage[margin=25mm]{geometry}
\usepackage{fontspec}
\setmainfont{lmroman10-regular.otf}[BoldFont=lmroman10-bold.otf,
 ItalicFont=lmroman10-italic.otf,BoldItalicFont=lmroman10-bolditalic.otf]
\usepackage{amsmath,amssymb,mathtools}
\usepackage{unicode-math}
\setmathfont{latinmodern-math.otf}
\AtBeginDocument{\let\setminus\smallsetminus}
\usepackage{xurl,hyperref}
\hypersetup{colorlinks=true,urlcolor=blue}
\setcounter{secnumdepth}{0}
\setlength{\parindent}{0pt}
\setlength{\parskip}{5pt}
\setlength{\emergencystretch}{3em}
\begin{document}
\section*{Kourovka Notebook Problem 21.23}\label{2532}
{\small UnsolvedMath ID 2532 · KOU-21.23 · Group Theory · Kourovka Notebook - New Problems, Issue 21}
\subsection{Definitions and mathematical statement}
For a finite group $G$, its enhanced power graph $E(G)$ is the finite
simple undirected graph whose vertices are the elements of $G$.
Distinct elements $x,y$ are adjacent exactly when the subgroup
$\langle x,y\rangle$ is cyclic, meaning generated by a single element.
The identity is included as a vertex. This adjacency condition is not
merely the condition that one of $x,y$ is a power of the other.

An induced cycle of length $n\ge4$ is a sequence of distinct vertices
$x_1,\ldots,x_n$ such that adjacent vertices in cyclic order are joined,
and no other pairs among these vertices are joined. In this graph,
consecutive pairs must generate cyclic groups while all nonconsecutive
pairs must generate noncyclic groups. A graph is chordal if it has no
induced cycle of any length at least four.

Part (a) asks, for each fixed integer $n\ge4$, for the complete set of
finite nonabelian simple groups $G$ for which $E(G)$ has no induced
$n$-cycle. Here simple means that $1$ and $G$ are its only normal
subgroups. Part (b) asks for the complete set of such groups for which
$E(G)$ is chordal.

A classification must give necessary and sufficient conditions, up to
group isomorphism. For part (a), excluding one specified length does
not exclude the others; part (b) imposes their simultaneous exclusion.
The auxiliary cograph notion in the source means absence of an induced
four-vertex path, a distinct restriction rather than either requested
classification.
\subsection{Short English statement}
Classify finite nonabelian simple groups whose enhanced power graphs avoid an induced cycle of a specified length, and those whose graphs are chordal.
\subsection{Sources}
\begin{itemize}
\item[S1] ULAM.AI and the UnsolvedMath Contributors, supplied problems.json, record 2532 (KOU-21.23), Kourovka Notebook - New Problems, Issue 21. Catalogue identity and source transcription; CC BY 4.0.
\url{https://huggingface.co/datasets/ulamai/UnsolvedMath}.
\item[S2] The Kourovka Notebook, 21st edition, new problems, Problem 21.23. Expanded formulation of the supplied catalogue statement.
\url{https://arxiv.org/pdf/1401.0300v47}.
\end{itemize}
\end{document}
